% Section fragment. Bibliography: GTOrbits2012, GTMobius2012.
\subsection{Quantitative factorization for centered weights}
\label{sec:gf-factorization}

We isolate the analytic implication needed after the arithmetic
progression main terms have been established. Throughout,
$L=\log(2N)$, all nilpotent groups are connected and simply connected,
and all filtrations are by closed connected rational subgroups. Constants
may depend on fixed arithmetic moduli and on the fixed dimension and
degree bounds.

\begin{proposition}[Factorization implication]\label{prop:gf-factorization}
Fix dimension and degree bounds $m,d$. Let $w_N:[N]\to\mathbb C$ satisfy
the following three hypotheses.
\begin{enumerate}
\item For a fixed $K_w\geq0$, $|w_N(n)|\ll L^{K_w}$.
\item For every $A_0,C_0>0$, uniformly over all interval progressions
$P\subseteq[N]$ of common difference at most $L^{C_0}$,
\begin{equation}\label{eq:gf-sw}
 \left|\sum_{n\in P}w_N(n)\right|
       \ll_{A_0,C_0}N L^{-A_0}.
\end{equation}
No lower bound on $|P|$ is imposed here.
\item There are fixed $c>0$, $C_1,K_1\geq0$ such that if a nilmanifold
of dimension at most $m$ and filtration degree at most $d$ has a
$Q$-rational adapted Mal'cev basis, $h$ is a polynomial sequence totally
$\delta$-equidistributed on $[N]$, and $F_0$ is Lipschitz with Haar mean
zero, then
\begin{equation}\label{eq:gf-ma}
 \left|\sum_{n\in P}w_N(n)F_0(h(n)\Gamma)\right|
 \ll N\delta^c Q^{C_1}L^{K_1}\|F_0\|_{\mathrm{Lip}}
\end{equation}
for every interval progression $P\subseteq[N]$ of modulus at most $Q$
and length at least $N/Q$. It is sufficient for this hypothesis to hold
only when
\begin{equation}\label{eq:gf-ma-regime}
 \delta\geq N^{-\sigma},\qquad Q\leq\delta^{-\kappa}
\end{equation}
for fixed positive $\sigma,\kappa$. An estimate allowing arbitrary
progression lengths is stronger than required.
\end{enumerate}
Then, for every $A,C_{\mathrm{init}}>0$, uniformly over initial
nilmanifolds of dimension at most $m$ and degree at most $d$ with an
$L^{C_{\mathrm{init}}}$-rational adapted basis, all polynomial sequences
$g$, and all Lipschitz $F$,
\begin{equation}\label{eq:gf-conclusion}
 \left|\sum_{n\leq N}w_N(n)F(g(n)\Gamma)\right|
     \ll_{A,C_{\mathrm{init}}}N L^{-A}\|F\|_{\mathrm{Lip}}.
\end{equation}
\end{proposition}

\begin{proof}
Normalize $\|F\|_{\mathrm{Lip}}=1$ and put
\[
 D=A+K_w+3,\qquad M_0=L^{\max\{1,C_{\mathrm{init}}\}}.
\]
Leave a factorization exponent $B$ unspecified for the moment. By
\cite[Theorem 1.19]{GTOrbits2012}, there is an integer $M$ with
\[
 M_0\leq M\ll_{B,m,d}M_0^{E(B,m,d)}
\]
and a factorization $g=\varepsilon g'\gamma$ into polynomial sequences.
Here $\varepsilon$ is $(M,N)$-smooth, $\gamma(n)\Gamma$ has period
$q\leq M$, and $g'$ is totally $M^{-B}$-equidistributed in a rational
subgroup quotient. The new adapted basis consists of $M$-rational
combinations of the initial one.

For each residue $j$ modulo $q$, choose a reduced rational representative
$\gamma_j$ of $\gamma(j)\Gamma$. Split that residue class into
$\lceil ML^D\rceil$ consecutive blocks of nearly equal cardinality.
For sufficiently large $N$ after $B$ is fixed, the resulting progressions
$P_{j,k}$ satisfy
\begin{equation}\label{eq:gf-cells}
 \#\{P_{j,k}\}\ll M^2L^D,\qquad
 |P_{j,k}|\gg\frac{N}{M^2L^D},\qquad
 \operatorname{diam}(P_{j,k})\ll\frac{N}{ML^D}.
\end{equation}
Their common difference is $q$. Indeed each original residue class has
cardinality comparable to $N/q$, and rounding the balanced block sizes
changes these estimates only by an absolute factor, since all powers of
$M$ are polylogarithmic in $N$ once $B$ is fixed.

Choose $n_{j,k}\in P_{j,k}$ and put $a_{j,k}=\varepsilon(n_{j,k})$.
The right-invariant ambient metric and smoothness imply
\begin{align*}
 &|F(\varepsilon(n)g'(n)\gamma_j\Gamma)
       -F(a_{j,k}g'(n)\gamma_j\Gamma)|\\
 &\hspace{1cm}\leq d(\varepsilon(n),\varepsilon(n_{j,k}))
       \ll L^{-D}\qquad(n\in P_{j,k}).
\end{align*}
There is no factor depending on the coordinates of $g'(n)$ in this
estimate: right invariance removes it. The total weighted freezing
error is therefore $O(NL^{K_w-D})=O(NL^{-A-3})$.

We use the quantitative conjugation claim proved in
\cite[p.~547]{GTMobius2012}, whose inputs are the rational-subgroup,
rational-element and metric comparisons of
\cite[Appendix A]{GTOrbits2012}. There are fixed $a>0$, $C_2\geq1$,
depending only on $m,d$, such that, with
\[
 H_j=\gamma_j^{-1}G'\gamma_j,\quad
 \Lambda_j=H_j\cap\Gamma,\quad
 h_j(n)=\gamma_j^{-1}g'(n)\gamma_j,
\]
the quotient $H_j/\Lambda_j$ has an adapted basis of rationality at most
$M^{C_2}$, $h_j$ is totally $M^{-aB+C_2}$-equidistributed, and
\[
 F_{j,k}(x\Lambda_j)=F(a_{j,k}\gamma_jx\Gamma)
 \quad\hbox{satisfies}\quad
 \|F_{j,k}\|_{\mathrm{Lip}}\leq M^{C_2},\qquad
 \|F_{j,k}\|_\infty\leq1.
\]
We have enlarged $C_2$ to absorb fixed multiplicative constants. These
exponents do not depend on $B$. In applying this published claim, one
must account for the commensurable lattice resulting from rational
conjugation; conjugation is not asserted to be an isometry of the
original quotients. Its finite-index and metric losses are included in
$a,C_2$.

Write $F_{j,k}=m_{j,k}+F_{j,k}^0$ by subtracting its Haar mean on
$H_j/\Lambda_j$. Then $|m_{j,k}|\leq1$,
$\int F_{j,k}^0=0$, and
$\|F_{j,k}^0\|_{\mathrm{Lip}}\leq2M^{C_2}$. This subtraction is needed
even if the initial function $F$ happened to have Haar mean zero.
For a dimension-zero factor the centered function simply vanishes.

Put
\[
 U=C_2+D+5,\qquad Q_*=M^U,\qquad b_*=aB-C_2.
\]
Since $M\geq L$, $Q_*$ dominates the rationality, cell modulus and
$N/|P_{j,k}|$. Applying \eqref{eq:gf-ma} to all centered cell functions
and using \eqref{eq:gf-cells} gives a total of at most
\begin{equation}\label{eq:gf-total-ma}
 C N L^{K_1}M^{-b_*c+C_1U+C_2+2+D}.
\end{equation}
Every cell is counted here. In particular, the bound $N$ in the
minor-arc hypothesis has not been replaced by an unsupported cell
cardinality.

We now fix $B$ sufficiently large that
\begin{equation}\label{eq:gf-choice-B}
 b_*>1,\qquad
 b_*c\geq C_1U+C_2+2+D+A+K_1+4,
 \qquad b_*\kappa\geq U+1
\end{equation}
when the restricted regime \eqref{eq:gf-ma-regime} is used. The last
condition is omitted for an unrestricted minor-arc estimate. These
inequalities involve only fixed exponents and not $E(B,m,d)$.
They make \eqref{eq:gf-total-ma} $O(NL^{-A-4})$. They also ensure
$Q_*\leq(M^{-b_*})^{-\kappa}$. The other regime condition,
$M^{-b_*}\geq N^{-\sigma}$, holds for all sufficiently large $N$:
after fixing $B$ the upper bound on $M$ is a fixed power of $\log N$.

Only at this point use that upper bound, writing $M\leq L^{E_*}$ after
enlarging the fixed exponent $E_*$. There are at most
$O(L^{2E_*+D})$ cells of modulus at most $L^{E_*}$. Apply
\eqref{eq:gf-sw} with
\[
 C_0=E_*+1,\qquad A_0=A+2E_*+D+4.
\]
Summing the constant terms $m_{j,k}\sum_{n\in P_{j,k}}w_N(n)$ gives
$O(NL^{-A-4})$. Combined with the centered estimates and freezing error,
this proves \eqref{eq:gf-conclusion}.
\end{proof}

\begin{remark}[Order of parameter choices]\label{rem:gf-order}
The order in the proof is
\[
 (A,K_w,C_{\mathrm{init}},c,C_1,K_1,a,C_2,\kappa,\sigma)
 \longrightarrow (D,U)\longrightarrow B\longrightarrow E_*
 \longrightarrow(A_0,C_0)\longrightarrow N_0.
\]
Thus there is no requirement that $B$ dominate the factorization upper
exponent $E(B,m,d)$. The centered progression estimate allows every
logarithmic saving, so its exponent can be chosen after $E_*$ is known.
The minor-arc input must have a fixed power form such as
\eqref{eq:gf-ma}; unspecified polylogarithmic specializations alone
would not justify this argument.
\end{remark}
